\documentclass[a4paper,14pt, unknownkeysallowed]{extreport}
\usepackage{placeins}
\usepackage{cmap} % Улучшенный поиск русских слов в полученном pdf-файле
\usepackage[T2A]{fontenc} % Поддержка русских букв
\usepackage[utf8]{inputenc} % Кодировка utf8
\usepackage[english,russian]{babel} % Языки: русский, английский
%\usepackage{pscyr} % Нормальные шрифты
\usepackage{enumitem}
\usepackage{multirow}
\usepackage[14pt]{extsizes}

\usepackage{caption}
\captionsetup{labelsep=endash}
\captionsetup[figure]{name={Рисунок}}

\usepackage{amsmath}

\usepackage{geometry}
\geometry{left=30mm}
\geometry{right=15mm}
\geometry{top=20mm}
\geometry{bottom=20mm}

\setlength{\parindent}{1.25cm} % для шрифта 14pt

\usepackage{titlesec}
\titleformat{\section}
	{\normalsize\bfseries}
	{\thesection}
	{1em}{}
\titlespacing*{\chapter}{0pt}{-30pt}{8pt}
\titlespacing*{\section}{\parindent}{*4}{*4}
\titlespacing*{\subsection}{\parindent}{*4}{*4}

\usepackage{setspace}
\onehalfspacing % Полуторный интервал

\frenchspacing
\usepackage{indentfirst} % Красная строка

\usepackage{titlesec}
\titleformat{\chapter}{\LARGE\bfseries}{\thechapter}{20pt}{\LARGE\bfseries}
\titleformat{\section}{\Large\bfseries}{\thesection}{20pt}{\Large\bfseries}

\usepackage{listings}
\usepackage{xcolor}

% Для листинга кода:
\lstset{ %
	language=lisp,   					% выбор языка для подсветки	
	basicstyle=\small\sffamily,			% размер и начертание шрифта для подсветки кода
	numbers=left,						% где поставить нумерацию строк (слева\справа)
	%numberstyle=,					% размер шрифта для номеров строк
	stepnumber=1,						% размер шага между двумя номерами строк
	numbersep=5pt,						% как далеко отстоят номера строк от подсвечиваемого кода
	frame=single,						% рисовать рамку вокруг кода
	tabsize=4,							% размер табуляции по умолчанию равен 4 пробелам
	captionpos=t,						% позиция заголовка вверху [t] или внизу [b]
	breaklines=true,					
	breakatwhitespace=true,				% переносить строки только если есть пробел
	escapeinside={\#*}{*)},				% если нужно добавить комментарии в коде
	backgroundcolor=\color{white},
}
\usepackage{pgfplots}
\usetikzlibrary{datavisualization}
\usetikzlibrary{datavisualization.formats.functions}

\usepackage{graphicx}
\newcommand{\img}[3] {
	\begin{figure}[h!]
		\center{\includegraphics[height=#1]{inc/img/#2}}
		\caption{#3}
		\label{img:#2}
	\end{figure}
}
\newcommand{\boximg}[3] {
	\begin{figure}[h]
		\center{\fbox{\includegraphics[height=#1]{inc/img/#2}}}
		\caption{#3}
		\label{img:#2}
	\end{figure}
}

% для матриц
\makeatletter
\renewcommand*\env@matrix[1][\arraystretch]{%
	\edef\arraystretch{#1}%
	\hskip -\arraycolsep
	\let\@ifnextchar\new@ifnextchar
	\array{*\c@MaxMatrixCols c}}
\makeatother

\usepackage{amsfonts}
\usepackage{graphicx}


\usepackage[justification=centering]{caption} % Настройка подписей float объектов

\usepackage[unicode,pdftex]{hyperref} % Ссылки в pdf
\hypersetup{hidelinks}

\usepackage{csvsimple}

\addto\captionsrussian{\def\refname{Список используемых источников}}

\newcommand{\code}[1]{\texttt{#1}}

\makeatletter
\renewcommand{\@biblabel}[1]{#1. \hfill}
\makeatother


\begin{document}


\begin{titlepage}
	\newgeometry{pdftex, left=2cm, right=2cm, top=2.5cm, bottom=2.5cm}
	\fontsize{12pt}{12pt}\selectfont
	\noindent \begin{minipage}{0.15\textwidth}
		\includegraphics[width=\linewidth]{inc/img/b_logo.jpg}
	\end{minipage}
	\noindent\begin{minipage}{0.9\textwidth}\centering
		\textbf{Министерство науки и высшего образования Российской Федерации}\\
		\textbf{Федеральное государственное бюджетное образовательное учреждение высшего образования}\\
		\textbf{«Московский государственный технический университет имени Н.Э.~Баумана}\\
		\textbf{(национальный исследовательский университет)»}\\
		\textbf{(МГТУ им. Н.Э.~Баумана)}
	\end{minipage}
	
	\noindent\rule{18cm}{3pt}
	\newline\newline
	\noindent ФАКУЛЬТЕТ $\underline{\text{«Информатика и системы управления»~~~~~~~~~~~~~~~~~~~~~~~~~~~~~~~~~~~~~~~~~~~~~~~~~~~~~~~}}$ \newline\newline
	\noindent КАФЕДРА $\underline{\text{«Программное обеспечение ЭВМ и информационные технологии»~~~~~~~~~~~~~~~~~~~~~~~}}$\newline\newline\newline\newline\newline\newline\newline
	
	
	\begin{center}
		\Large\textbf{Отчет по лабораторной работе №2}
		\Large\textbf {по дисциплине <<Функциональное и логическое программирование>>}
	\end{center}
	
	\noindent\textbf{Тема} $\underline{\text{~~Определений функций пользователя~~~~~~~~~~~~~~~~~~~~~~~~~~~~~~~~~~~~~~~~~~~~~~~~~~~~~~~~~~~~~~~~~}}$\newline\newline\newline
	\noindent\textbf{Студент} $\underline{\text{~~Котляров Н.А.~~~~~~~~~~~~~~~~~~~~~~~~~~~~~~~~~~~~~~~~~~~~~~~~~~~~~~~~~~~~~~~~~~~~~~~~~~~~~~~~~~~~~~~~~}}$\newline\newline
	\noindent\textbf{Группа} $\underline{\text{~~ИУ7-61Б~~~~~~~~~~~~~~~~~~~~~~~~~~~~~~~~~~~~~~~~~~~~~~~~~~~~~~~~~~~~~~~~~~~~~~~~~~~~~~~~~~~~~~~~~~~~~~~~~}}$\newline\newline
	\noindent\textbf{Оценка (баллы)} $\underline{\text{~~~~~~~~~~~~~~~~~~~~~~~~~~~~~~~~~~~~~~~~~~~~~~~~~~~~~~~~~~~~~~~~~~~~~~~~~~~~~~~~~~~~~~~~~~~~~~~~~~~}}$\newline\newline
	\noindent\textbf{Преподаватель} $\underline{\text{~~Толпинская Н.~Б.~~Строганов~Ю.~В.~~~~~~~~~~~~~~~~~~~~~~~~~~~~~~~~~~~~~~~~~~~~~~~~~~}}$\newline
	
	\begin{center}
		\vfill
		Москва~---~\the\year
		~г.
	\end{center}
	\restoregeometry
\end{titlepage}



\section*{Задание 1}

\begin{enumerate}
	\item Найти длину списка (по верхнему уровню);
	\item Найти сумму элементов числового списка;
	\item Найти сумму элементов числового списка, стоящих на нечетных позициях исходного
	списка (нумерация от 0);
	\item Сформировать список из элементов числового списка, больших заданного значения;
	\item Удалить заданный элемент из списка (один или все вхождения).
	\item Объединить два списка.
\end{enumerate}


\begin{lstlisting}
my_length(List, Length) :- my_length(List, 0, Length).
my_length([], Acc, Acc) :- !.
my_length([_|T], Acc, Length) :- !,
	NewAcc is Acc + 1,
	my_length(T, NewAcc, Length).

my_sum_list(List, Sum) :- my_sum_list(List, 0, Sum), !.
my_sum_list([], Acc, Acc) :- !.
my_sum_list([H|T], Acc, Sum) :- 
	NewAcc is Acc + H,
	my_sum_list(T, NewAcc, Sum).

my_odd_sum(List, Sum) :- !,
	my_odd_sum(List, 0, Sum).
my_odd_sum([], Sum, Sum) :- !.
my_odd_sum([_], Sum, Sum).
my_odd_sum([_, X|T], Acc, Sum) :-
	NewAcc is Acc + X,
	my_odd_sum(T, NewAcc, Sum).

new_big_list([], _, []) :- !.
new_big_list([X|Xs], Y, [X|Ys]) :- X > Y, !, new_big_list(Xs, Y, Ys).
new_big_list([X|Xs], Y, Ys) :- !, new_big_list(Xs, Y, Ys).

del_member(_, [], []) :- !.
del_member(X, [X|Tail], Tail) :- !.
del_member(X, [Y|Tail], [Y|Tail1]) :-
	del_member(X, Tail, Tail1).

del_all_member(_, [], []) :- !.
del_all_member(X, [X|Tail], Result) :- !,
	del_all_member(X, Tail, Result).
del_all_member(X, [Y|Tail], [Y|Tail1]) :- del_all_member(X, Tail, Tail1).

concratinate_lists([], L, L) :- !.
concratinate_lists([H|T], L, [H|Result]) :- concratinate_lists(T, L, Result).
\end{lstlisting}



\begin{table}[]
	\centering
	\caption{Задание 1}
	\label{task1}
	\begin{tabular}{|p{0.8cm}|p{4cm}|p{7.5cm}|p{4cm}|}
		\hline
		\textbf{N} & \textbf{Состояние резольвенты} & \textbf{Для каких термов запускается алгоритм унификации и результат подстановки} & \textbf{Дальнейшие действия} \\ 
		\hline
		1 & my\_length([1,2,3], Res). & my\_length([1,2,3], 0). Результат: Унификация прошла успешно, переменная Acc принимает значение 0. & Прямой ход. Создается новое состояние резольвенты с целью my\_length([1, 2,3], Acc, Res), Res = Length, Acc = 0. \\
		\hline
		2 & my\_length([1,2,3], Acc, Res), Res = Length, Acc = 0. & my\_length([1,2,3], 0). Результат: Унификация прошла успешно, переменная Acc принимает значение 0. & Прямой ход. Создается новое состояние резольвенты с целью my\_length([1, 2,3], Acc, Res), Res = Length, Acc = 0. \\
		\hline
		2 & my\_length([2,3], 1, Res). & my\_length([2,3], 1, Length). Результат: Унификация прошла успешно, переменная Acc принимает значение 1. & Прямой ход. Создается новое состояние резольвенты с целью my\_length([3], 2, Res). \\
		\hline
		3 & my\_length(List, Acc, Res), Res=0, Acc=1, List=[2,3]. &my\_lengthh(List, Acc, Res)my\_lengthth([\_|T],Acc, Length). Результат: Унификация прошла успешно, переменная NewAcc принимает значение Acc + 1 = 2, List=T. & Прямой ход. Создается новое состояние резольвенты с целью my\_length(List, Acc, Rew), Res = Res, Acc=NewAcc, List=T. \\
		\hline
		
	\end{tabular}
\end{table}


\begin{table}[]
	\centering
	\caption{Задание 1}
	\label{task1}
	\begin{tabular}{|p{0.8cm}|p{4cm}|p{7.5cm}|p{4cm}|}
		\hline
		\textbf{N} & \textbf{Состояние резольвенты} & \textbf{Для каких термов запускается алгоритм унификации и результат подстановки} & \textbf{Дальнейшие действия} \\ 
		\hline
		4 & my\_length(List, Acc, Res), Res=0, Acc=2, List=[3]. &my\_lengthh(List, Acc, Res)my\_lengthth([\_|T],Acc, Length). Результат: Унификация прошла успешно, переменная NewAcc принимает значение Acc + 1 = 3, List=T. & Прямой ход. Создается новое состояние резольвенты с целью my\_length(List, Acc, Rew), Res = Res, Acc=NewAcc, List=T. \\
		\hline
		5 & my\_length(List, Acc, Res), Res=0, Acc=3, List=[] &my\_lengthh([], Acc, Acc). Результат: Унификация прошла успешно, переменная Acc принимает значение 3. & Обратный ход. Возвращаемся к состоянию 4, где выполняется подстановка: Res=3. \\
		\hline
		6 & my\_length(List, Acc, Res), Res=3, Acc=3, List=[] &my\_lengthh([], Acc, Acc). Результат: Унификация прошла успешно. & Обратный ход. Возвращаемся к состоянию 3, где выполняется подстановка: Res=3. \\
		\hline
		7 & my\_length([1,2,3], Res). &my\_lengthh([1,2,3], 0). Результат: Унификация прошла успешно. & Обратный ход. Возвращаемся к состоянию 2, где выполняется подстановка: Length=3. \\
		\hline
	\end{tabular}
\end{table}

\end{document}